\subsection{Generalized Continum Hypotesis (GCH)}
An additional axiom that can be added to ZF is the Generalized Continum Hypotesis (GCH). It 
states that for every cardinality $\kappa$, its successor is $2^{\kappa}$. This is what we 
will assume from now on. We will show how this axiom can be used to prove that every set 
can be well ordered (e.g. it can be put in bijection with some ordinal).

\subsection{Sierpinski's theorem}
The Sierpinski's theorem states that the GCH implies that for every infinite set $A$ the cardinality of the cartesian product of $A$ 
with itself is the same as the cardinality of $A$ itself. This means that $A \times A \sim A$ for every infinite set $A$.

\begin{lemma}
    Let $k$ be a cardinality of an infinite set, then $k + 1 = k$ under the assumption that the generalized continum 
    hypotesys holds true at least for $k$ (local GCH for $k$).
\end{lemma}
\begin{proof}
    $$ k \leq k + 1 < 2^k                                                  $$
    $$ \texttt{GCH(k)}: \exists \bar{k} : k \leq \bar{k} < 2^k \Rightarrow \bar{k} = k $$
    $$ \therefore k + 1 = k                                                $$
\end{proof}

\begin{lemma}
    Let $k$ be a cardinality of an infinite set, then $2k = k$ under the assumption that the generalized continum 
    hypotesys holds true at least for $k$ (local GCH for $k$).
\end{lemma}
\begin{proof}
    $$ k \leq k + k < 2^k + 2^k = 2^{k+1} = 2^k                                        $$
    $$ \texttt{GCH(k)}: \exists \bar{k} : k \leq \bar{k} < 2^k \Rightarrow \bar{k} = k $$
    $$ \therefore 2k = k                                                               $$
\end{proof}

\begin{theorem}[Sierpinski's theorem]
    Let $k$ be a cardinality of an infinite set, then $k^2 = k$ under the assumption that the generalized continum 
    hypotesys holds true at least for $k$.
\end{theorem}
\begin{proof}
    $$ k \leq k^2 < (2^k)^2 = 2^{2k} = 2^k                                             $$
    $$ \texttt{GCH(k)}: \exists \bar{k} : k \leq \bar{k} < 2^k \Rightarrow \bar{k} = k $$
    $$ \therefore k^2 = k                                                              $$
\end{proof}

Notice how we can assume that GCH holds true in general and then prove that for every infinite set $A$ with cardinality 
$k$, since $k = k^2$ then $A \times A \sim A$. This clearly holds true for $\mathbb{N}$, but 
with the GCH this also applies to $\mathbb{R}$ and to every infinite set in general. \newpage